\documentclass[11pt,reqno]{amsart}

%==============================================================
% A Numerical Variance Budget for Iterated Brownian Motion
%==============================================================

\usepackage[utf8]{inputenc}
\usepackage[T1]{fontenc}
\usepackage{lmodern}
\usepackage[margin=1.15in]{geometry}
\usepackage{amsmath,amssymb,amsthm,mathtools}
\usepackage{microtype}
\usepackage{graphicx}
\usepackage{booktabs}
\usepackage{hyperref}
\usepackage{cleveref}

% Pull figures from the project's outputs/ tree at build time.
\graphicspath{{../outputs/figures/}}

\theoremstyle{plain}
\newtheorem{theorem}{Theorem}[section]
\newtheorem{proposition}[theorem]{Proposition}
\newtheorem{lemma}[theorem]{Lemma}
\newtheorem{corollary}[theorem]{Corollary}

\theoremstyle{definition}
\newtheorem{definition}[theorem]{Definition}

\theoremstyle{remark}
\newtheorem{remark}[theorem]{Remark}

\newcommand{\R}{\mathbb{R}}
\newcommand{\N}{\mathbb{N}}
\newcommand{\E}{\mathbb{E}}
\newcommand{\Var}{\mathrm{Var}}
\newcommand{\Cov}{\mathrm{Cov}}
\newcommand{\Proj}{\mathrm{Proj}}
\newcommand{\HX}[1]{\mathcal{H}^{X}_{#1}}
\newcommand{\HY}[1]{\mathcal{H}^{Y}_{#1}}
\newcommand{\Loc}{L}
\newcommand{\IBM}{Z}

\title{A Numerical Variance Budget\\
       for Iterated Brownian Motion}
\author{Ramiro Fontes}
\date{\today}

\begin{document}

\begin{abstract}
Subtracting two analytic chaos projections from the naive Monte
Carlo estimator of $\E\!\int_0^T \varphi(X_{Y_s})\,ds$ for $\varphi$
odd at $T = 1$ removes $62\%$ of the variance and yields a
$2.6\times$ effective-sample-size multiplier, at essentially no
runtime cost. The projections come from a fully analytic
$3 \times 3$ source-resolved variance budget for the iterated
Brownian motion (IBM) occupation functional
$F_T(\varphi) = \int_0^T \varphi(\IBM_s)\,ds$
over bidegrees $(m, n) \in \{0, 1, 2\}^2$, evaluated in closed form
using the source projection formula of the companion paper together
with the $n = 2$ clock chaos of the Brownian local time, derived
here from Tanaka's formula and matching the Nualart--Vives chaos
expansion of $\Loc_T^x(Y)$. Across three test functions ($\varphi$
even, odd, mixed) and three time horizons ($T = 0.25, 1, 4$) the
parity selection rules of the companion paper hold to machine
precision, the $3 \times 3$ truncation captures one half to over
four-fifths of $\Var(F_T)$, and across a $9 \times 6$ sweep of test
functions, horizons, and control sets, predicted and realised
reductions agree closely; a split-sample check gives the same
conclusion. A finite-dimensional Hermite tensor visualisation
(Stage 2) is exploratory and lives in an appendix.
\end{abstract}

\maketitle

%==============================================================
\section{Introduction}
\label{sec:intro}
%==============================================================

The companion paper introduces a source-resolved chaos expansion for
functionals of iterated Brownian motion (IBM), $\IBM_t = X_{Y_t}$,
where $X$ is a two-sided Brownian scenery and $Y$ is an independent
Brownian clock. Working in the product Wiener chaos
$L^2(\Omega) = \bigoplus_{m,n \ge 0} \HX{m} \otimes \HY{n}$, that
paper proves a source projection formula for the canonical occupation
functional
\[
  F_T(\varphi)
   \;=\; \int_0^T \varphi(\IBM_s)\,ds
   \;=\; \int_\R \varphi(X_x)\,\Loc_T^x(Y)\,dx,
   \qquad \varphi \in C_c^\infty(\R),
\]
and derives parity selection rules together with explicit closed-form
projections at low bidegree.

This is its numerical companion. We compute, for three concrete
test functions and three time horizons, the variance share of every
analytic projection $\Proj_{m,n} F_T(\varphi)$ in the truncated set
$(m, n) \in \{0, 1, 2\}^2$. The full $3 \times 3$ truncation
(including the previously uncomputed cells $(0, 2), (1, 2), (2, 2)$)
is a fully analytic accounting: every cell is a closed-form spatial
integral, and every parity zero is exact up to numerical roundoff.

\Cref{fig:headline} (Section~\ref{sec:stage1}) shows the $3 \times 3$
analytic variance share for the nine $(\varphi, T)$ combinations.
The parity selection rules of the companion paper are visually
obvious: for $\varphi$ even, the column $m = 1$ is zero and so is
the cell $(0, 1)$; for $\varphi$ odd, the rows $m = 0$ and $m = 2$
are zero throughout; for $\varphi$ mixed all cells are active except
$(0, 1)$, which cancels because $\E[\varphi_{\mathrm{mixed}}(X_x)]$
is even in $x$ while $k_T(x, s)$ is odd in $x$, so the spatial
integral defining $\Proj_{0, 1} F_T$ cancels (Section~\ref{sec:setup}).
The captured share of $\Var(F_T)$ ranges from $47\%$ to $84\%$ across
the nine cases (Table~\ref{tab:stage1}), with the residual living in
higher bidegrees that we do not compute analytically here.

The main application is a sparse, source-resolved control variate.
For $\varphi$ odd at $T = 1$, the analytic projections
$\Proj_{1,0} F_T$ and $\Proj_{1,1} F_T$ together carry $59\%$ of
$\Var(F_T)$. Subtracting them per sample gives an unbiased Monte
Carlo estimator with $62\%$ empirical variance reduction (matching
the Stage 1 prediction within sampling noise) and an effective
sample size $2.6\times$ that of the naive estimator
(Section~\ref{sec:control-variate}).

The variance-accounting layer (Stage 1) uses closed-form projections
evaluated per Monte Carlo sample, with sample-variance estimates
over 4000 realisations. A separate visualisation layer (Stage 2)
regresses $F_T$ onto an orthonormal Hermite-tensor basis built from
finite indicator factors, extending the bidegree decomposition to
the wider plane $(m, n) \in \{0, \dots, 5\}^2$ at the cost of a
basis-coverage residual that is large when the analytic projection
kernels live in spatial directions the indicator basis cannot
resolve. Stage 1 is the primary accounting; Stage 2 is exploratory
and is relegated to an appendix.

Chaos-expansion methods for local-time functionals also appear in
the work of Hu and Nualart on renormalised self-intersection local
time for fractional Brownian motion \cite{HuNualart}; the present
computations differ by resolving two independent Gaussian sources,
the scenery $X$ and the clock $Y$, rather than encoding chaos order
in a single source. We do not derive a new PDE; we do not push
Stage 2 to large $K$ beyond the appendix; we do not claim asymptotic
theorems for the $T \to 0$ or $T \to \infty$ limits of the variance
budget. Where we note empirical scaling patterns we say so
explicitly.

%==============================================================
\section{Setup and analytic background}
\label{sec:setup}
%==============================================================

Let $X$ be a two-sided standard Brownian motion (the scenery) and $Y$
a standard Brownian motion on $[0, T]$ independent of $X$ (the
clock). Set $\IBM_t = X_{Y_t}$. For $\varphi \in C_c^\infty(\R)$, the
occupation functional is
\[
  F_T(\varphi)
   \;=\; \int_0^T \varphi(X_{Y_s})\,ds
   \;=\; \int_\R \varphi(X_x)\,\Loc_T^x(Y)\,dx,
\]
with $\Loc_T^x(Y)$ the local time of $Y$ at level $x$ up to $T$.

The companion paper proves the source projection formula
\begin{equation}
  \Proj_{m, n} F_T(\varphi)
   \;=\; \int_\R \Proj^X_m \varphi(X_x)
        \,\cdot\, \Proj^Y_n \Loc_T^x(Y)\,dx,
   \qquad m, n \ge 0,
   \label{eq:source-projection}
\end{equation}
in $L^2(\Omega)$, with the integral converging as a Bochner integral
in $\HX{m} \otimes \HY{n}$. We use \eqref{eq:source-projection} as
the analytic accounting backbone for the entire numerical budget.

The scenery factors are Hermite polynomials in $X_x \sim N(0, |x|)$:
\begin{equation}
  \Proj^X_0 \varphi(X_x) = \E[\varphi(X_x)],\quad
  \Proj^X_1 \varphi(X_x) = \E[\varphi'(X_x)]\,X_x,\quad
  \Proj^X_2 \varphi(X_x) = \tfrac{1}{2}\E[\varphi''(X_x)](X_x^2 - |x|).
  \label{eq:scenery-proj}
\end{equation}
The clock projections at $n = 0, 1$ are
\begin{align}
  \Proj^Y_0 \Loc_T^x(Y)
   &= p_T(x) \;=\; \int_0^T \frac{1}{\sqrt{2 \pi s}}\,
     e^{-x^2/(2s)}\,ds,\\
  \Proj^Y_1 \Loc_T^x(Y)
   &= \int_0^T k_T(x, s)\,dY_s,\quad
   k_T(x, s) = \mathrm{erf}\!\bigl(\tfrac{x}{\sqrt{2s}}\bigr)
              - \mathrm{erf}\!\bigl(\tfrac{x}{\sqrt{2T}}\bigr).
\end{align}

\subsection{The clock projection at $n = 2$}
\label{sec:proj2-clock}
The clock projection at $n = 2$ is computed here from Tanaka's
formula. With $\psi_t(x) = (2 \pi t)^{-1/2} e^{-x^2/(2t)}$ the
Brownian heat kernel, projecting each piece of the Tanaka identity
onto $\HY{2}$ gives:

\begin{lemma}[$\Proj^Y_2 \Loc_T^x(Y)$]
\label{lem:proj2-LT}
For $x \in \R$ and $T > 0$,
\begin{equation}
  \Proj^Y_2 \Loc_T^x(Y)
   \;=\; \psi_T(x)\,\bigl(Y_T^2 - T\bigr)
        \,-\, 2 \int_0^T \psi_s(x)\,Y_s\,dY_s,
   \label{eq:proj2-LT}
\end{equation}
equivalently $I_2^Y\bigl(K_T(\cdot, \cdot; x)\bigr)$ as a double
Wiener--It\^o integral with symmetric kernel
$K_T(s_1, s_2; x) = \psi_T(x) - \psi_{\max(s_1, s_2)}(x)$. The kernel
is in $L^2([0, T]^2)$ for each $x$, including $x = 0$. In any case,
changing the value at the singleton $\{x = 0\}$ does not affect the
spatial Bochner integral of \eqref{eq:source-projection}. This is the
$n = 2$ specialisation of the Nualart--Vives chaos expansion of
Brownian local time \cite{NualartVives}, with the Brownian-heat-kernel
normalisation $\psi_t(x) = (2 \pi t)^{-1/2} e^{-x^2/(2t)}$ used here.
\end{lemma}

\begin{proof}
By Tanaka,
$\Loc_T^x(Y) = |Y_T - x| - |x| - \int_0^T \mathrm{sgn}(Y_s - x)\,dY_s$.
We project each piece onto $\HY{2}$ separately.

\emph{(i) The endpoint term.}
$|Y_T - x|$ is a function of the single Gaussian $Y_T \sim N(0, T)$.
Its Hermite expansion at variance $T$ reads
$|Y_T - x| = \sum_{n \ge 0} (n!)^{-1}\,
 \E\!\bigl[(\partial_y)^n |Y_T - x|\bigr]\, h_n(Y_T; T)$,
with $h_n(z; T)$ the Hermite polynomial of variance $T$
($h_2(z; T) = z^2 - T$). The second derivative
$(\partial_y)^2 |y - x| = 2 \delta(y - x)$ gives
$\E[(\partial_y)^2 |Y_T - x|] = 2 \psi_T(x)$, and hence
\[
  \Proj^Y_2 |Y_T - x|
   \;=\; \tfrac{1}{2!}\,(2 \psi_T(x))\,(Y_T^2 - T)
   \;=\; \psi_T(x)(Y_T^2 - T)
   \;=\; \psi_T(x)\,I_2^Y(\mathbf{1}_{[0,T]} \otimes \mathbf{1}_{[0,T]}).
\]

\emph{(ii) The deterministic term.} $|x|$ is constant in $\omega$ and
its $\HY{2}$ projection is $0$.

\emph{(iii) The It\^o-integral term.} Write $u_s = \mathrm{sgn}(Y_s - x)$;
$u$ is adapted with $\E[u_s] = -\mathrm{erf}(x/\sqrt{2s})$
($\HY{0}$ part) and first-chaos part
$\Proj^Y_1 u_s = (1!)^{-1}\,\E[u'(Y_s)]\,Y_s = 2 \psi_s(x)\,Y_s$
(again using $\E[\mathrm{sgn}'(Y_s - x)] = 2 \psi_s(x)$). For an
adapted process $u_s$ with chaos expansion
$u_s = \sum_m I_m^Y(g_m^{(s)})$ on $\{s_1 \le \dots \le s_m \le s\}$,
the It\^o integral $\int_0^T u_s\,dY_s$ has chaos at order $m+1$,
so the second-chaos contribution comes from the first-chaos part of
$u$ alone:
\[
  \Proj^Y_2 \int_0^T \mathrm{sgn}(Y_s - x)\,dY_s
   \;=\; \int_0^T \Proj^Y_1 u_s\,dY_s
   \;=\; \int_0^T 2 \psi_s(x)\,Y_s\,dY_s.
\]
Writing $Y_s = \int_0^s dY_{s_1}$ and converting the iterated It\^o
integral to a symmetric multiple Wiener integral,
\[
  \int_0^T 2 \psi_s(x)\,Y_s\,dY_s
   \;=\; 2 \int_0^T \!\int_0^s \psi_s(x)\,dY_{s_1}\,dY_s
   \;=\; I_2^Y\!\bigl(\psi_{\max(s_1, s_2)}(x)\bigr),
\]
the symmetrisation following from the convention
$I_2^Y(f) = 2! \int_{0 < s_1 < s_2 < T} f(s_1, s_2)\,dY_{s_1}\,dY_{s_2}$
applied to a symmetric $f$.

\emph{(iv) Combining.}
Subtracting (iii) from (i) and using (ii),
\[
  \Proj^Y_2 \Loc_T^x(Y)
   \;=\; \psi_T(x)\,I_2^Y(\mathbf{1} \otimes \mathbf{1})
        \,-\, I_2^Y(\psi_{\max(s_1, s_2)}(x))
   \;=\; I_2^Y\!\bigl(\psi_T(x) - \psi_{\max(s_1, s_2)}(x)\bigr),
\]
which is \eqref{eq:proj2-LT} in kernel form.

\emph{(v) $L^2$ check.}
By symmetry of $K_T$ and the change of variables $u = \max(s_1, s_2)$,
\[
  \|K_T(\cdot, \cdot; x)\|_{L^2([0,T]^2)}^2
   \;=\; 2 \int_0^T s\,\bigl(\psi_T(x) - \psi_s(x)\bigr)^2\,ds.
\]
Each factor of $\psi_t(x) - \psi_T(x)$ is bounded uniformly in
$t \in (0, T]$ for $x \ne 0$ (since $\psi_t(x)$ decays
super-polynomially as $t \to 0^+$ for $x \ne 0$). At $x = 0$, the
integrand $s\,\psi_s(0)^2 = s\,/(2\pi s) = 1/(2\pi)$ is bounded as
$s \to 0^+$: the $\psi_s(0)^2 \sim 1/(2\pi s)$ singularity is
cancelled by the area density $s\,ds$ for
$\max(s_1, s_2) = s$. Hence $K_T(\cdot, \cdot; x) \in L^2([0, T]^2)$
for every $x \in \R$.
\end{proof}

\noindent\emph{Normalisation cross-check.} Computing
$\Var(\Proj^Y_2 \Loc_T^x(Y))$ two ways, once as $2 \|K_T\|_{L^2}^2$
and once as the orthogonal decomposition
$\Var(\psi_T(x)(Y_T^2 - T)) + \Var(2 \int \psi_s Y_s\,dY_s)
- 2\,\Cov(\cdot, \cdot)$, yields the same closed form
$2 T^2 \psi_T(x)^2 + 4 \int_0^T s\,\psi_s(x)^2\,ds
- 8 \psi_T(x) \int_0^T s \psi_s(x)\,ds$, fixing the normalisation.
The implementation reads the formula \eqref{eq:proj2-LT} directly via
a left-point It\^o discretisation.

\paragraph{Parity selection rules.}
For $\varphi$ even, $\E[\varphi^{(m)}(X_x)] = 0$ for odd $m$; hence
$\Proj_{m, n} F_T = 0$ for every odd $m$ and any $n$. Additionally
the $n$-th clock-chaos kernel of $\Loc_T^x(Y)$ has parity
$(-1)^n$ in $x$, while $\E[\varphi(X_x)]$ is even in $x$ for
$\varphi$ even, so $\Proj_{0, n} F_T = 0$ for odd $n$. For $\varphi$
odd, $\E[\varphi^{(m)}(X_x)] = 0$ for even $m$; hence
$\Proj_{m, n} F_T = 0$ for every even $m$ and any $n$. For $\varphi$
mixed, no parity rule applies. We summarise:

\begin{center}
\small
\begin{tabular}{c|ccc}
$\varphi$ even & $n=0$ & $n=1$ & $n=2$ \\\hline
$m=0$ & nonzero & \textbf{zero} & nonzero \\
$m=1$ & \textbf{zero} & \textbf{zero} & \textbf{zero} \\
$m=2$ & nonzero & nonzero & nonzero \\
\end{tabular}
\hspace{0.7cm}
\begin{tabular}{c|ccc}
$\varphi$ odd & $n=0$ & $n=1$ & $n=2$ \\\hline
$m=0$ & \textbf{zero} & \textbf{zero} & \textbf{zero} \\
$m=1$ & nonzero & nonzero & nonzero \\
$m=2$ & \textbf{zero} & \textbf{zero} & \textbf{zero} \\
\end{tabular}
\end{center}

%==============================================================
\section{Stage 1: the analytic 3\texttimes 3 variance budget}
\label{sec:stage1}
%==============================================================

We evaluate \eqref{eq:source-projection} sample by sample on a fine
grid for the nine bidegrees $(m, n) \in \{0, 1, 2\}^2$, and aggregate
sample variances over Monte Carlo realisations. The spatial grid is
symmetric on $[-A, A]$ with $A = 5$ and $N_x = 400$ intervals, the
time grid is $0 = t_0 < t_1 < \dots < t_{N_t} = T$ with $N_t = 400$
steps, and the two-sided scenery is sampled by independent Gaussian
increments outward from $X_0 = 0$. Per sample, we compute on the
spatial grid:
\begin{itemize}
  \item $V_1(x) = \int_0^T k_T(x, s)\,dY_s$ via left-point It\^o
    discretisation;
  \item $V_2(x) = \Proj^Y_2 \Loc_T^x(Y)$ via the closed form of
    \cref{lem:proj2-LT};
  \item conditional Gaussian moments $\E[\varphi^{(j)}(X_x)]$,
    $j = 0, 1, 2$, by 40-point Gauss--Hermite quadrature in the
    standardised variable.
\end{itemize}
The nine $\Proj_{m,n} F_T$ are formed by trapezoidal integration in
$x$ of the appropriate products, and sample variances use
$N_1 = 4000$ Monte Carlo realisations. The three test functions are
\[
  \varphi_{\mathrm{even}}(x) = e^{-x^2},\quad
  \varphi_{\mathrm{odd}}(x) = x e^{-x^2},\quad
  \varphi_{\mathrm{mixed}}(x) = (x + x^2)\,e^{-x^2},
\]
and the three horizons are $T \in \{0.25, 1, 4\}$.
\Cref{tab:stage1} reports each cell's variance share of
$\Var(F_T(\varphi))$ across the nine $(\varphi, T)$ cases.

\begin{table}[h]
\centering
\small
\begin{tabular}{lcccccccccc}
\toprule
$\varphi$, $T$ &
$(0,0)$ & $(1,0)$ & $(0,1)$ & $(1,1)$ & $(2,0)$ & $(2,1)$ &
$(0,2)$ & $(1,2)$ & $(2,2)$ & $\Sigma$\\
\midrule
even, 0.25  & --- & 0     & 0     & 0     & 0.283 & 0.249 & 0.264 & 0     & 0.048 & 0.844 \\
even, 1.0   & --- & 0     & 0     & 0     & 0.280 & 0.204 & 0.205 & 0     & 0.034 & 0.723 \\
even, 4.0   & --- & 0     & 0     & 0     & 0.287 & 0.149 & 0.180 & 0     & 0.026 & 0.642 \\[2pt]
odd, 0.25   & --- & 0.454 & 0     & 0.246 & 0     & 0     & 0     & 0.025 & 0     & 0.726 \\
odd, 1.0    & --- & 0.442 & 0     & 0.182 & 0     & 0     & 0     & 0.023 & 0     & 0.646 \\
odd, 4.0    & --- & 0.382 & 0     & 0.118 & 0     & 0     & 0     & 0.024 & 0     & 0.524 \\[2pt]
mixed, 0.25 & --- & 0.427 & 0     & 0.214 & 0.031 & 0.013 & 0.017 & 0.027 & 0.002 & 0.731 \\
mixed, 1.0  & --- & 0.374 & 0     & 0.151 & 0.017 & 0.003 & 0.009 & 0.020 & 0.003 & 0.577 \\
mixed, 4.0  & --- & 0.318 & 0     & 0.110 & 0.009 & 0.006 & 0.005 & 0.020 & 0.006 & 0.474 \\
\bottomrule
\end{tabular}
\caption{Stage 1 variance shares $\Var(\Proj_{m,n} F_T) / \Var(F_T)$
for the analytic $3 \times 3$ truncation. The $(0,0)$ column carries
$\E[F_T]^2$ not variance and is omitted. Entries shown as $0$ are
predicted-zero by parity (Section~\ref{sec:setup}) and come out
numerically at the level of $10^{-30}$ to $10^{-40}$ (numerical roundoff).
$\Sigma$ is the total share captured by the $3\times 3$ truncation;
the residual $1 - \Sigma$ lives in higher bidegrees.}
\label{tab:stage1}
\end{table}

\Cref{fig:headline} shows a $3 \times 3$ grid of
3$\times$3 heatmaps displaying the entries of \cref{tab:stage1} as
filled cells. The parity zero pattern is visually obvious in each
panel. By the interaction-energy identity in the companion paper,
the $(1, 1)$ entries of \cref{tab:stage1} are empirical estimates of
$\|K_T\|_{L^2(\R \times [0, T])}^2 / \Var(F_T)$, i.e.\ the normalised
first scenery--clock interaction energy.

\begin{figure}[t]
\centering
\includegraphics[width=0.99\textwidth]{stage1_heatmap_3x3.png}
\caption{Stage 1 analytic variance share by bidegree $(m, n) \in \{0, 1, 2\}^2$
for each of the nine $(\varphi, T)$ cases. Cells labelled \texttt{cst}
are the $(0, 0)$ constant component and carry $\E[F_T]^2$, not variance.
Cells labelled \texttt{0} are predicted-zero by parity selection rules
(see Section~\ref{sec:setup}); empirically these entries come out at
machine precision, i.e.\ $10^{-30}$--$10^{-40}$, as the scenery
expectations they contain are zero exactly by Gaussian symmetry before
any Monte Carlo step. Coloured cells carry the analytic mass: pure
clock $(0,2)$ and pure scenery $(2,0)$ for $\varphi$ even, the column
$m = 1$ for $\varphi$ odd, all cells for $\varphi$ mixed except
$(0, 1)$.}
\label{fig:headline}
\end{figure}

\paragraph{Empirical orthogonality.}
The companion paper proves orthogonality of distinct bidegree
components in $L^2(\Omega)$. Across all nine cases the empirical
correlation matrix among the nine projection arrays has diagonal $1$
and off-diagonal entries at the level of Monte Carlo noise:
$\max |\rho_{ij}| \le 0.10$, against a theoretical standard error
$1/\sqrt{N_1} \approx 0.016$ per off-diagonal entry. The slight
excess over pure sampling-noise heuristics reflects finite-grid
discretisation; orthogonality is exact for the continuum chaos
components, while the reported components are evaluated on the
finite simulation grid. Per-case correlation heatmaps are produced
as part of the pipeline and saved alongside the other figures.

\subsection{$T$-dependence: the shrinking mixed share}
\label{sec:T-dep}
At fixed $\varphi$, the relative weight of mixed-bidegree cells
versus pure-scenery cells \emph{shrinks} as $T$ grows. Concretely,
for $\varphi$ odd the ratio $\Var(\Proj_{1, 1} F_T) / \Var(\Proj_{1, 0} F_T)$
goes from $0.54$ at $T = 0.25$ to $0.41$ at $T = 1$ to $0.31$
at $T = 4$. For $\varphi$ mixed the same ratio falls from $0.50$ to
$0.40$ to $0.34$. For $\varphi$ even, the cell $(2, 1)$ as a fraction
of $(2, 0)$ falls from $0.88$ to $0.73$ to $0.52$. The pure-scenery
sector loses ground only mildly with $T$, while mixed and pure-clock
shares fall faster. We do not attempt to extract an asymptotic
theorem from this; we record it as an observed empirical pattern in
the parameter range studied. It is consistent with the heuristic
that as $T$ grows, the mean local time $p_T(x)$ widens and the
spatial weighting that drives $\Proj^Y_n$ at $n \ge 1$ smooths out,
shifting variance to higher clock bidegrees beyond the $3 \times 3$
truncation.

\medskip
\noindent\emph{Heuristic context.} For small $T$, the clock $Y$ stays
near $0$ and $F_T$ becomes nearly deterministic after scaling, so
the budget concentrates in low-order bidegrees and the
$3 \times 3$ truncation captures most of $\Var(F_T)$. For larger
$T$, $Y$ explores a wider range, higher bidegrees enter, and a
fixed $3 \times 3$ truncation captures less. This frames the
captured-share trend $\Sigma_{T = 0.25} > \Sigma_{T = 1} > \Sigma_{T = 4}$
visible in \cref{tab:stage1} as observation, not theorem.

\subsection{Robustness checks}
\label{sec:robustness}
We record three standalone validations of the analytic budget.

\paragraph{Direct $\|K_T\|^2$ cross-check.}
The companion paper's interaction-energy identity
$\Var(\Proj_{1, 1} F_T) = \|K_T\|^2$ predicts the $(1, 1)$ cell
variance analytically. Using
$\langle \sigma_x, \sigma_y \rangle_{L^2(\R)}
 = \min(|x|, |y|)\,\mathbf{1}_{xy \ge 0}$
the kernel norm reduces to a deterministic 3D quadrature
\[
  \|K_T\|^2 = 2\!\int_0^A\!\!\int_0^A \E[\varphi'(X_u)]\,\E[\varphi'(X_v)]
  \,\min(u, v)
  \int_0^T k_T(u, s)\,k_T(v, s)\,ds\,du\,dv.
\]
For $\varphi_{\mathrm{odd}}$ and $\varphi_{\mathrm{mixed}}$ at
$T = 1$, the quadrature value matches the empirical Stage 1 $(1, 1)$
share of \cref{tab:stage1} within $|0.004|$ absolute share
(\cref{tab:kernel-norm-check}), confirming the interaction-energy
identity numerically and independently of any Monte Carlo estimate
of $\Var(\Proj_{1,1} F_T)$.

\begin{table}[h]
\small\centering
\begin{tabular}{lrrrrr}
\toprule
$\varphi$ & $T$ & $\|K_T\|^2$ &
$\widehat\Var(\Proj_{1,1} F_T)$ &
analytic share & empirical share \\
\midrule
odd   & 1.0 & 0.00654 & 0.00668 & 0.178 & 0.182 \\
mixed & 1.0 & 0.00654 & 0.00645 & 0.153 & 0.151 \\
\bottomrule
\end{tabular}
\caption{Direct analytic $\|K_T\|^2$ cross-check vs Stage 1
empirical $\widehat\Var(\Proj_{1,1} F_T)$ for the $(1, 1)$ cell.
Shares are normalised by the case-specific $\widehat\Var(F_T)$.
$\|K_T\|^2$ depends on $\varphi$ only through $\E[\varphi'(X_x)]$,
which is identical for $\varphi_{\mathrm{odd}}$ and
$\varphi_{\mathrm{mixed}}$ (their $\phi'$ have the same even part).}
\label{tab:kernel-norm-check}
\end{table}

\paragraph{Discretisation sensitivity.}
We sweep the spatial half-width $A \in \{3, 5, 10\}$ and the grid
sizes $(N_x, N_t) \in \{(200, 200), (400, 400), (800, 800)\}$ for
$\varphi_{\mathrm{mixed}}$ at $T = 1$ at $N_1 = 2000$ MC samples per
setting (9 settings, total runtime $\le 5\,\mathrm{s}$). Across the
9 settings the captured share $\Sigma$ of the 3$\times$3 truncation
lies in $[0.560, 0.620]$, the $(1, 0)$ share in $[0.350, 0.401]$,
the $(1, 1)$ share in $[0.140, 0.175]$, and the $(0, 2)$ share in
$[0.0089, 0.0132]$. The variation is comparable to the MC noise
band $\sim 1/\sqrt{N_1} \approx 2\%$ and shows no monotone drift
under refinement; the budget is not a discretisation artefact.

\paragraph{Additivity.}
For $\varphi_{\mathrm{sum}} = \varphi_{\mathrm{even}} + \varphi_{\mathrm{odd}}$,
the parity rules of \cref{sec:setup} predict cell-by-cell additivity
$\Var(\Proj_{m, n} F_T(\varphi_{\mathrm{sum}}))
 = \Var(\Proj_{m, n} F_T(\varphi_{\mathrm{even}}))
 + \Var(\Proj_{m, n} F_T(\varphi_{\mathrm{odd}}))$,
since exactly one of the two summands contributes to each cell by
parity (the cross covariance is zero in $L^2(\Omega)$). Running
Stage 1 on $\varphi_{\mathrm{sum}}$ with the same parameters and
comparing entry-by-entry against the corresponding rows of
\cref{tab:stage1}, we find across the seven non-zero cells and
three horizons RMS relative deviation $\le 0.08$ and maximum
relative deviation $\le 0.19$ on any single cell, consistent with
independent-MC noise of two runs at $N_1 = 4000$ each. The bidegree
budget respects the parity decomposition of the observable.

%==============================================================
\section{Source-resolved control variate}
\label{sec:control-variate}
%==============================================================

For any subset $\mathcal{A} \subset \{(m, n) \in \N^2 : (m, n) \ne (0, 0)\}$
of non-constant bidegrees, define the source-resolved control
variate
\begin{equation}
  C_\mathcal{A} \;\coloneqq\; \sum_{(m, n) \in \mathcal{A}} \Proj_{m, n} F_T(\varphi).
  \label{eq:CA-def}
\end{equation}
Each summand is a non-constant chaos element, so $\E[C_\mathcal{A}] = 0$
and the controlled estimator $F_T - C_\mathcal{A}$ has the same mean
$\E[F_T]$ as the naive estimator $F_T$. By orthogonality of distinct
bidegrees in $L^2(\Omega)$,
\begin{equation}
  \Var\!\bigl(F_T - C_\mathcal{A}\bigr)
   \;=\; \Var(F_T)
        \,-\, \sum_{(m, n) \in \mathcal{A}} \Var\!\bigl(\Proj_{m, n} F_T\bigr).
  \label{eq:CA-var}
\end{equation}
The variance reduction realised by $C_\mathcal{A}$ is therefore exactly
the sum of Stage 1 variance shares over $\mathcal{A}$.

\subsection{The 62\% reduction at $(\varphi_{\mathrm{odd}}, T = 1)$}
\label{sec:cv-demo}
For $(\varphi, T) = (\varphi_{\mathrm{odd}}, 1)$ and
$\mathcal{A} = \{(1, 0), (1, 1)\}$ the projections $\Proj_{1, 0} F_T$
and $\Proj_{1, 1} F_T$ are deterministic spatial integrals against a
single linear functional of $(X, Y)$ already constructed to evaluate
$F_T$, so the runtime overhead is essentially free. Run with
$N = 8000$ Monte Carlo realisations:

\begin{table}[h]
\centering
\small
\begin{tabular}{lr}
\toprule
quantity & value \\
\midrule
naive sample variance $\widehat\Var(F_T)$        & $0.0372$ \\
controlled sample variance $\widehat\Var(F_T - C_\mathcal{A})$  & $0.0141$ \\
empirical variance reduction                       & $62.1\%$ \\
predicted reduction (unrounded Stage 1 shares)     & $58.7\%$ \\
effective sample size multiplier                   & $2.6\times$ \\
total wall-clock time                              & $0.80\,\mathrm{s}$ \\
\bottomrule
\end{tabular}
\caption{Source-resolved control variate for
$(\varphi, T) = (\varphi_{\mathrm{odd}}, 1)$,
$\mathcal{A} = \{(1, 0), (1, 1)\}$, $N = 8000$. The empirical
reduction matches the predicted reduction within sampling noise.
The predicted value uses unrounded Stage 1 shares; reading the
rounded entries from \cref{tab:stage1} (a separate 4000-sample run,
$0.442 + 0.182 = 0.624$) gives a similar magnitude.}
\label{tab:cv}
\end{table}

\noindent The controlled estimator with $N = 8000$ gives the same
precision as a naive estimator with $N \approx 21{,}000$. The
$62.1\%$ empirical reduction and the $58.7\%$ predicted reduction
agree up to Monte Carlo noise; the finite-sample provenance of the
small gap is discussed in the caveat below.

\subsection{Validation across the full sweep}
\label{sec:control-validation}
The single-case result above is one entry in a broader validation
that the source-resolved variance
budget predicts realised control-variate variance reduction across all
nine $(\varphi, T)$ cases. For each case we evaluate six control sets
$\mathcal{A}$:
\begin{itemize}\itemsep0pt
  \item \textbf{top1 / top2 / top3:} the one, two, three largest
    non-constant Stage 1 cells by variance share;
  \item \textbf{full\_3x3:} all non-constant cells in
    $\{0, 1, 2\}^2$ with non-zero share;
  \item \textbf{total\_order\_1:} $\{(1, 0), (0, 1)\}$ (cells with
    $m + n = 1$);
  \item \textbf{total\_order\_2:} $\{(1, 0), (0, 1), (1, 1), (2, 0), (0, 2)\}$
    (cells with $1 \le m + n \le 2$).
\end{itemize}
For each $(\varphi, T, \mathcal{A})$ we record the predicted reduction
$\sum_{(m,n) \in \mathcal{A}}\Var(\Proj_{m,n} F_T) / \Var(F_T)$ and the
empirical reduction $1 - \widehat\Var(F_T - C_\mathcal{A}) / \widehat\Var(F_T)$,
both from the same $N = 8000$ Monte Carlo run; in this experiment we
use the spatial-form $F_T = \int_\R \varphi(X_x)\,L_T^x(Y)\,dx$ (with
$L_T^x$ via Tanaka) so that the discrete $F_T$ and the discrete
projections share the same spatial discretisation, reducing a small
finite-grid mismatch present in the path-form $F_T$ used elsewhere in
the paper.

\Cref{fig:budget-validation} (left) shows the resulting scatter of
empirical against predicted reduction across the $9 \times 6 = 54$
trials. All points lie close to the diagonal $y = x$: the RMS error
is $0.018$ and the maximum absolute error is $0.049$. The budget
predicts realised reduction across the full sweep; the small spread
about the diagonal is consistent with finite-$N$ Monte Carlo noise
plus minor finite-grid effects.

\Cref{fig:budget-validation} (right) compares source-aware top-$k$
controls against blind total-order controls on the same axes. For
$\varphi$ even, $\{(1, 0), (0, 1)\}$ is the entire total-order-$1$
set but both cells are predicted-zero by parity; the corresponding
bar is at zero in the figure (the light-grey bar is invisible).
The source-aware top-$2$ control selects two analytically-non-zero
cells (here $(2, 0)$ and $(0, 2)$) and achieves roughly $50\%$
variance reduction with the same number of components. Across all
nine cases, source-aware top-$k$ at small $k$ matches or beats
total-order at much larger cardinality: e.g. top-$2$ ($k=2$ cells)
reaches the same reduction as total-order-$2$ ($\le 5$ cells, of
which several are parity-zero) for both $\varphi$ even and $\varphi$
odd.

\begin{figure}[h]
\centering
\includegraphics[width=0.49\textwidth]{predicted_vs_empirical_reduction.png}\hfill
\includegraphics[width=0.49\textwidth]{source_aware_vs_total_order.png}
\caption{\emph{Left:} predicted vs empirical control-variate
variance reduction across the nine $(\varphi, T)$ cases and six
control sets ($N = 8000$); the diagonal is $y = x$. RMS deviation
$0.018$, maximum $0.049$. \emph{Right:} empirical reduction for
source-aware top-$k$ vs total-order control sets. For $\varphi$ even,
total-order-$1$ has both cells predicted-zero by parity and the bar
is at $0$; source-aware top-$2$ picks the largest non-zero cells
instead.}
\label{fig:budget-validation}
\end{figure}

\noindent\emph{Reading.} The source-resolved budget is operational:
it predicts realised variance reduction from analytic chaos
controls. Source-aware sparse controls avoid structurally zero cells
and can match or outperform total-order controls of the same or
greater cardinality.

\medskip
\noindent\emph{Split-sample robustness.} To rule out same-sample
artefacts, we also ran a split-sample validation in which the
budget and the selected cells are estimated on one Monte Carlo
batch (sample $A$, $N_A = 8000$) and the controlled estimator is
evaluated on an \emph{independent} batch (sample $B$, $N_B = 8000$,
disjoint seed). The predicted-vs-realised reductions remain close
to the diagonal: across the same nine $(\varphi, T)$ cases and six
control sets, RMS error $0.015$, maximum absolute error $0.022$,
all within the same Monte Carlo noise band as the same-sample
check. The split-sample scatter is recorded alongside the other
figures in the project outputs.

\subsection{Comparison with the conditional Monte Carlo baseline}
\label{sec:cond-baseline}
The classical conditional Monte Carlo estimator
$\widehat F^{\mathrm{cond}}_T = F_T - \bigl(\E[F_T \mid Y] - \E[F_T]\bigr)$
removes all $Y$-randomness from $F_T$ that is captured by
$\E[F_T \mid Y]$. By the chaos decomposition,
$\E[F_T \mid Y] - \E[F_T] = \sum_{n \ge 1} \Proj_{0, n} F_T$,
i.e.\ the entire pure-clock sector. The conditioning baseline is
therefore a special case of the source-resolved control restricted
to the column $m = 0$.

For $\varphi_{\mathrm{odd}}$, $\E[\varphi(X_x)] \equiv 0$ so the
pure-clock sector is zero and the conditioning baseline gives no
reduction; this is one source of the parity selection rule. We
compare on $\varphi_{\mathrm{mixed}}$ and $\varphi_{\mathrm{even}}$
at $T = 1$ (\cref{tab:cond-comparison}, $N = 8000$). For
$\varphi_{\mathrm{mixed}}$ the bulk of $\Var(F_T)$ lives in $m \ge 1$
cells, so conditioning yields only a $1.8\%$ reduction; the
source-resolved top-2 control selects $\{(1, 0), (1, 1)\}$ and
removes $55\%$. For $\varphi_{\mathrm{even}}$ conditioning captures
the $m = 0$ sector and reaches $25\%$; source-resolved top-3 reaches
$66\%$. A combined estimator subtracting both the conditioning term
and the Stage 1 $m \ge 1$ cells (avoiding double-counting the
$m = 0$ sector) achieves the best reduction in each row. The
takeaway is not a runtime benchmark: source-resolved sparse
controls dominate the naive conditioning baseline whenever the
variance has substantial $m \ge 1$ content.

\begin{table}[h]
\small\centering
\begin{tabular}{llrrrl}
\toprule
$\varphi$ & estimator & $\widehat\Var$ & reduction & ESS factor & details \\
\midrule
mixed & naive          & 0.0416 & --     & $1.0\times$  & $F_T$ \\
mixed & conditioning   & 0.0408 & $1.8\%$ & $1.02\times$ & $F_T - C_{\mathrm{cond}}$ \\
mixed & top-2          & 0.0188 & $55\%$  & $2.21\times$ & $\{(1,0), (1,1)\}$ \\
mixed & top-3          & 0.0180 & $57\%$  & $2.31\times$ & $\{(1,0), (1,1), (1,2)\}$ \\
mixed & combined       & 0.0161 & $61\%$  & $2.58\times$ & $C_{\mathrm{cond}} +$ $m \ge 1$ cells \\[2pt]
even  & naive          & 0.0380 & --     & $1.0\times$  & $F_T$ \\
even  & conditioning   & 0.0283 & $25\%$  & $1.34\times$ & $F_T - C_{\mathrm{cond}}$ \\
even  & top-2          & 0.0197 & $48\%$  & $1.93\times$ & $\{(2,0), (0,2)\}$ \\
even  & top-3          & 0.0128 & $66\%$  & $2.98\times$ & $\{(2,0), (0,2), (2,1)\}$ \\
even  & combined       & 0.0095 & $75\%$  & $4.01\times$ & $C_{\mathrm{cond}} +$ $m \ge 1$ cells \\
\bottomrule
\end{tabular}
\caption{Source-resolved controls vs the conditioning baseline at
$T = 1$, $N = 8000$. The ``combined'' row subtracts both
$C_{\mathrm{cond}}$ and the Stage 1 cells with $m \ge 1$ (not the
full Stage 1 top-$k$ as the latter would double-count the $m = 0$
sector already in $C_{\mathrm{cond}}$).}
\label{tab:cond-comparison}
\end{table}

\paragraph{Caveat.}
This is not a Monte Carlo benchmark. We fix the control coefficient
at $1$ on each selected component. This is the population
regression-optimal coefficient for exact chaos controls: if
$C = \Proj_{m, n} F_T$ with $(m, n) \ne (0, 0)$, then by orthogonality
of distinct bidegrees in $L^2(\Omega)$,
\[
  \Cov(F_T, C)
   \;=\; \Cov(\Proj_{m, n} F_T,\,\Proj_{m, n} F_T)
   \;=\; \Var(C),
\]
so the standard optimal control-variate coefficient
$\beta^\star = \Cov(F_T, C) / \Var(C)$ equals $1$. For a sum of
selected bidegrees $C_\mathcal{A} = \sum_{(m, n) \in \mathcal{A}}
\Proj_{m, n} F_T$, the same calculation separates componentwise: the
cross-bidegree covariances are zero, $\Cov(F_T, C_\mathcal{A}) =
\sum_{(m, n) \in \mathcal{A}} \Var(\Proj_{m, n} F_T) =
\Var(C_\mathcal{A})$, and again $\beta^\star = 1$.

At finite grid and finite $N$, empirical orthogonality is only
approximate (Section~\ref{sec:stage1}), so fitted coefficients could
reduce sample variance slightly further. We keep $\beta = 1$ to
preserve the analytic, unbiased, non-fitted control-variate
construction. The same caveats apply to the broader validation of
\cref{sec:control-validation}: this is not a Monte Carlo benchmark,
we do not optimise coefficients empirically, finite-grid and
finite-$N$ effects can cause small deviations from the diagonal of
\cref{fig:budget-validation}, and we do not compare against
importance-sampling or quasi-Monte Carlo alternatives nor study how
the analytic projections depend on the spatial half-width $A$ or
grid sizes $N_x, N_t$ as these are refined. The purpose is to verify
that the analytic variance budget identifies effective control
components, not to claim a production-grade variance-reduction
recipe.

%==============================================================
\section{Conclusion}
\label{sec:conclusion}
%==============================================================

The full $3 \times 3$ analytic variance budget is computable in
closed form via the source projection formula plus the explicit
$\Proj^Y_n \Loc_T^x(Y)$ kernels for $n = 0, 1, 2$. For every test
case, parity selection rules predicted by the companion paper
hold to machine precision, and the captured share of $\Var(F_T)$ is
substantial: $47$--$84\%$ of the total. As a concrete application,
subtracting the $(1, 0)$ and $(1, 1)$ projections from the naive
Monte Carlo estimator of $\E[F_T(\varphi_{\mathrm{odd}})]$ at
$T = 1$ removes $62\%$ of the variance, matching the Stage 1
prediction of $59\%$ within sampling noise.

We see two natural follow-ups. First, the mixed-vs-pure shrinkage
pattern in Section~\ref{sec:T-dep} is an empirical observation
across three time horizons; whether it admits a clean asymptotic
formulation as $T \to \infty$ requires more work. Second, a fully
optimised control-variate estimator with regression-tuned
coefficients, larger $(m, n)$ truncation, and benchmarking against
existing variance-reduction baselines is out of scope here and is
the natural subject of a follow-up paper.

\subsection*{Reproducibility}
The Python pipeline that produces every figure and CSV is provided
alongside the paper in the same repository. A single end-to-end
driver script runs Stage 1, the appendix Stage 2 sweep and the
control-variate experiment in a few minutes on a recent laptop; the
control-variate cell alone runs in under a second.

\appendix

%==============================================================
\section{Stage 2: finite-dimensional Hermite tensor visualisation}
\label{sec:stage2}
%==============================================================

This appendix records an exploratory visualisation that projects
$F_T(\varphi)$ onto an orthonormal Hermite tensor basis built from
$K_x = K_y = K$ piecewise-constant indicator factors per source. The
visualisation extends to bidegrees $(m, n) \in \{0, \dots, 5\}^2$ and
is useful for showing the qualitative shape of the chaos budget at
$(m, n)$ that Stage 1 does not compute analytically. The cost is a
basis-coverage residual that can be large.

We project paths of $X$ and $Y$ onto $K_x$ and $K_y$ orthonormal
piecewise-constant indicator basis functions in $L^2$, yielding
factors $\zeta^X \in \R^{K_x}$ and $\zeta^Y \in \R^{K_y}$ that are
approximately iid $N(0, 1)$. We standardise each factor in sample so
that the empirical second moments are exactly $1$. The orthonormal
Hermite tensor basis is
$\widetilde H_{\alpha, \beta}(\zeta^X, \zeta^Y)
 = \prod_j h_{\alpha_j}(\zeta^X_j)/\sqrt{\alpha_j!}
 \cdot \prod_k h_{\beta_k}(\zeta^Y_k)/\sqrt{\beta_k!}$,
with $h_m$ the probabilist's Hermite polynomial. We sweep multi-indices
$\alpha \in \N^{K_x}, \beta \in \N^{K_y}$ with total degree
$|\alpha| \le 5$ and $|\beta| \le 5$.

For each $(\alpha, \beta)$ we estimate $c_{\alpha, \beta}^2$ by the
unbiased U-statistic
$\widehat{c_{\alpha, \beta}^2} =
\bigl[(\sum_s F_T^{(s)} \widetilde H_{\alpha, \beta}^{(s)})^2
 - \sum_s (F_T^{(s)} \widetilde H_{\alpha, \beta}^{(s)})^2\bigr]
/ (N (N - 1))$,
which can be negative on cells where the true value is zero. The
share of bidegree $(m, n)$ is
$\max(\sum_{|\alpha| = m, |\beta| = n} \widehat{c_{\alpha, \beta}^2},\, 0)
/ \Var(F_T)$, with clipping applied only at the bidegree-aggregated
level for visualisation. The raw signed values are retained in the
Stage 2 output table for downstream diagnostics.

At $K = 3$, residual share is $60$--$90\%$; at $K = 7$ the U-statistic
saturates for $\varphi$ even at low $\Var(F_T)$ because the feature
count $\binom{K + 5}{5}^2$ approaches the sample size. We do not
draw quantitative conclusions from saturated entries. The Stage 2
heatmaps qualitatively respect the parity rules of
Section~\ref{sec:setup} and reproduce the Stage 1 dominant cells
where the basis can resolve them; some cells, including $(1, 0)$ for
$\varphi$ odd and $(2, 1)$ for $\varphi$ even, are systematically
under-attributed by Stage 2. These discrepancies are basis-coverage
effects of the finite indicator subspace; Stage 1 gives the rigorous
source attribution. The full set of heatmaps, refinement
diagnostics and an auto-generated summary report ship with the
pipeline. For example, for $\varphi$ even at $T = 4$ the residual
outside the $3 \times 3$ block is concentrated mainly in cells such
as $(0, 4)$, $(4, 1)$, $(2, 3)$, consistent with higher even
bidegrees rather than leakage into forbidden parity cells; visible
Stage 2 mass in odd-$m$ cells (e.g. $(3, 2)$, $(1, 5)$) is the
basis-coverage / U-statistic noise artefact noted above, not a
violation of the parity selection rule.

\Cref{fig:stage2-even} shows a representative Stage 2 heatmap for
$\varphi$ even at $T = 1$, displaying the $(m, n) \in \{0, \dots, 5\}^2$
plane that Stage 1 does not analytically reach. The dominant Stage 2
cell at $(0, 2)$ matches the Stage 1 cell $(0, 2)$ in
\cref{tab:stage1}, and the parity-zero column $m = 1$ together with
the parity-zero cell $(0, 1)$ are visibly empty (modulo small
Monte-Carlo noise leakage); higher even-$m$ rows show additional
mass that Stage 1's $3 \times 3$ truncation does not cover.

\begin{figure}[h]
\centering
\includegraphics[width=0.55\textwidth]{heatmap_even_T1.0.png}
\caption{Representative Stage 2 heatmap for $\varphi$ even at
$T = 1$, $K_x = K_y = 3$, $\mathrm{deg\_max} = 5$, $N_2 = 10000$.
Cells are coloured by Stage-2 share of $\Var(F_T)$. The dominant
cell $(0, 2)$ matches the Stage 1 entry; column $m = 1$ and cell
$(0, 1)$ are empirically empty, in agreement with the parity rule.
Higher-$(m, n)$ cells visible here lie outside the $3 \times 3$
truncation of Stage 1.}
\label{fig:stage2-even}
\end{figure}

\begin{thebibliography}{9}

\bibitem{AllouZheng}
H.~Allouba and W.~Zheng,
\emph{Brownian-time processes: the PDE connection and the half-derivative generator},
Ann. Probab. \textbf{29} (2001), no.~4, 1780--1795.

\bibitem{ChaosPaper}
R.~Fontes,
\emph{Source-Resolved Chaos Expansions for Iterated Brownian Motion},
companion paper, 2026.

\bibitem{Janson}
S.~Janson,
\emph{Gaussian Hilbert Spaces},
Cambridge Tracts in Mathematics, vol.~129,
Cambridge University Press, Cambridge, 1997.

\bibitem{KhoshnevisanLewis}
D.~Khoshnevisan and T.~M.~Lewis,
\emph{Stochastic calculus for Brownian motion on a Brownian fracture},
Ann. Appl. Probab. \textbf{9} (1999), no.~3, 629--667.

\bibitem{Lee}
A.~J.~Lee,
\emph{U-Statistics: Theory and Practice},
Statistics: Textbooks and Monographs, vol.~110,
Marcel Dekker, New York, 1990.

\bibitem{Nualart}
D.~Nualart,
\emph{The Malliavin Calculus and Related Topics},
2nd ed., Probability and Its Applications, Springer-Verlag, Berlin, 2006.

\bibitem{NualartVives}
D.~Nualart and J.~Vives,
\emph{Chaos expansions and local times},
Publ. Mat. \textbf{36} (1992), no.~2B, 827--836.

\bibitem{HuNualart}
Y.~Hu and D.~Nualart,
\emph{Renormalized self-intersection local time for fractional Brownian motion},
Ann. Probab. \textbf{33} (2005), no.~3, 948--983.

\bibitem{RevuzYor}
D.~Revuz and M.~Yor,
\emph{Continuous Martingales and Brownian Motion},
3rd ed., Grundlehren der Mathematischen Wissenschaften, vol.~293,
Springer-Verlag, Berlin, 1999.

\end{thebibliography}

\end{document}
